\section{Specification}
\label{sec:spec}

We present the library specification following standard software
documentation conventions: interface first
(\S\ref{sec:interface}), public API
(\S\ref{sec:core-spec}--\S\ref{sec:app-spec}), internal lemmas
(\S\ref{sec:internal}).

\subsection{Interface: types and structures}
\label{sec:interface}

Throughout, \texttt{K} denotes a topological space (the parameter
space) and \texttt{f : K $\to$ $\mathbb{R}$} a continuous scalar
energy.

In the proof of Theorem~\ref{thm:AP-combinatorial}, the
Lemma~3.1 of~\cite{DLT13} supplies, for each near-critical
parameter $t$, a local replacement on an
$\mathcal{H}^n$-disjoint pair of open sets in $M^{n+1}$ whose
boundary energy is reduced by at least $1/(4N)$ on a
neighborhood of $t$. The replacement family
$\{\partial \Omega_{t,\tau}\}_{\tau \in [0,1]}$ at a
near-critical $t$ has its Hausdorff measure decreased by at
least $1/(4N)$ inside $U_t'$. We keep only the scalar
\emph{replacement energy} $E : K \to \mathbb{R}$ and the
uniform saving inequality; the geometric family is discarded.

\begin{definition}[\texttt{LocalWitness} structure]
\label{def:witness}
\leanmargin{LocalWitness}{Witness.lean}
A \texttt{LocalWitness K f t $\varepsilon$} carries an open
neighborhood $U \ni t$ in $K$, a continuous replacement energy
$E : K \to \mathbb{R}$, and the inequality
$\forall s \in U,\; f(s) - E(s) \ge \varepsilon$.
\end{definition}

The proof of the Claim works on the set $K \subset [0,1]$ of
parameters at which the sweepout is within $1/N$ of the min-max
value. We use the same convention.

\begin{definition}[Near-critical set]
\label{def:NC}
\leanmargin{nearCritical}{OneDim/InitialCover.lean}
For continuous $f : K \to \mathbb{R}$, scalar $m_0$, and
$\delta \ge 0$, the $\delta$-near-critical set is
$\{t \in K : f(t) \ge m_0 - \delta\}$. When $\delta = 1/N$ we
write $\NC$.
\end{definition}

At the end of the refinement the data is a family
$\{\overline{J_l}\}$ covering $K$, with each $J_l$ carrying a
replacement whose energy saving is at least $1/(4N)$ and with
two-fold overlap. We abstract this output: the near-critical
threshold $1/N$ becomes a parameter $\delta$ and the saving floor
$1/(4N)$ a parameter $\varepsilon \le \delta$.

\begin{definition}[\texttt{FiniteCoverWithWitnesses} structure]
\label{def:ref}
\leanmargin{FiniteCoverWithWitnesses}{Cover.lean}
A \texttt{FiniteCoverWithWitnesses K f m\textsubscript{0} $\delta$ $\varepsilon$}
carries a finite nonempty index type $\iota$, pieces
$\{p_l \subseteq K\}_{l \in \iota}$, replacement energies
$\{E_l : K \to \mathbb{R}\}$, and savings
$\{s_l \in \mathbb{R}_{>0}\}$, subject to four conditions:
(I) $f - E_l \ge s_l$ on $p_l$;
(II) each $t \in K$ lies in at most two pieces;
(III) $s_l \ge \varepsilon$ for all $l$;
(IV) $\{t \in K : f(t) \ge m_0 - \delta\} \subseteq \bigcup_l p_l$.
\end{definition}

The skip-2 spacing condition (a) of the original proof is a
purely one-dimensional geometric condition on a family of open
intervals. We isolate it in its own structure so the
disjointness lemmas that depend only on spacing can be proved
independently of the witness data.

\begin{definition}[\texttt{SkippedSpacedIntervals} structure]
\label{def:spaced}
\leanmargin{SkippedSpacedIntervals}{OneDim/SpacedIntervals.lean}
A \texttt{SkippedSpacedIntervals} records a natural number $n$,
centers $c_1, \ldots, c_n \in \unit$, positive radii
$r_1, \ldots, r_n > 0$, and the skip-2 spacing axiom
$c_i + r_i < c_{i+2} - r_{i+2}$ whenever $i+2 \le n$. The
$i$-th open interval $I_i = (c_i - r_i, c_i + r_i) \cap \unit$
is an abbreviation on this structure.
\end{definition}

On top of the geometric structure, the initial cover also
carries a witness center $w_i \in \NC$ with an attached local
witness for each interval, together with the containment of
each $I_i$ inside the corresponding witness neighborhood. In our
formalization the interval center $c_i$ and the witness center
$w_i$ are recorded separately: grid points need not themselves
lie in $\NC$.

\begin{definition}[\texttt{InitialCover} structure]
\label{def:ic}
\leanmargin{InitialCover}{OneDim/InitialCover.lean}
Specialized to $K = \unit$. An
\texttt{InitialCover f m\textsubscript{0} N} carries the data of
a \texttt{SkippedSpacedIntervals} (Definition~\ref{def:spaced},
satisfying spacing condition (a)), plus, for each
$i = 1, \ldots, n$, a witness center $w_i \in \NC$ and a local
witness at $w_i$ with saving $1/(4N)$, subject to (b)~the
containment $I_i \subseteq U_{w_i}$ and (c) the coverage
$\NC \subseteq \bigcup_i I_i$.
\end{definition}

The refinement is described as an iteration: each $I_j$ with
$j \ge 2$ is either kept whole or split into two halves. The
state of this iteration after $\ell$ moves is the following.

\begin{definition}[\texttt{PartialRefinement} structure]
\label{def:pr}
\leanmargin{PartialRefinement}{OneDim/PartialRefinement.lean}
Given \texttt{ic : InitialCover} and $\ell \in \mathbb{N}$, a
\texttt{PartialRefinement ic $\ell$} carries pieces
$J_1, \ldots, J_\ell \subseteq \unit$ and an index assignment
$\sigma : \{1, \ldots, \ell\} \to \{1, \ldots, n\}$ with
(i) $J_k \subseteq I_{\sigma(k)}$ for each $k$;
(ii) $I_{\sigma(k)} \subseteq \bigcup_{k'} J_{k'}$ for each $k$.
\end{definition}

\subsection{Public API: combinatorial core}
\label{sec:core-spec}

The library's main theorem is the existence of a refined
multiplicity-bounded cover of the near-critical set carrying the
witness savings: the non-trivial combinatorial content extracted
from the original proof. Its hypotheses coincide with the intro
statement (Theorem~\ref{thm:intro-scalar}); the output is a
\texttt{FiniteCoverWithWitnesses} (Definition~\ref{def:ref}).

\begin{theorem}[Combinatorial core]
\label{thm:core}
\leanmargin{exists\_refinement}{OneDim/Assembly.lean}
Let $f : \unit \to \mathbb{R}$ be continuous with
$m_0 = \sup f$, and let $N \in \mathbb{N}$ with $N > 0$.
Suppose every $t \in \NC$ admits a \texttt{LocalWitness} at $t$
with saving $1/(4N)$ (Definition~\ref{def:witness}). Then there
exists a
\texttt{FiniteCoverWithWitnesses}$(\unit, f, m_0, 1/N, 1/(4N))$
(Definition~\ref{def:ref}): a finite family of closed pieces
$\{p_l \subseteq \unit\}$, continuous replacement energies
$\{E_l : \unit \to \mathbb{R}\}$, and savings
$\{s_l \ge 1/(4N)\}$ satisfying conditions~(I)--(IV).
\end{theorem}

The proof is given in~\S\ref{sec:proof} as three chained stages:
Stage~A builds an initial cover from a Lebesgue-number argument
on $\NC$; Stage~B runs a bounded smallest-index induction to a
terminal partial refinement; Stage~C assembles the output by
extending the open-neighborhood saving bound to the closure via
continuity and deriving two-fold overlap via the parity-gap
disjointness argument.

\subsection{Public API: sup-reduction corollaries}
\label{sec:app-spec}

The arithmetic that turns a
\texttt{Finite\discretionary{-}{}{}Cover\discretionary{-}{}{}With\discretionary{-}{}{}Witnesses}
into a competitor with reduced supremum is a short scalar lemma
independent of covers, pieces, or multiplicity: any competitor
that saves a fixed amount on the relevant super-level set yields
the same amount of supremum reduction.

\begin{lemma}[Scalar sup reduction]
\label{lem:scalar-sup-reduction}
\leanmargin{csSup\_range\_le\_of\_pointwise\_saving}{Cover.lean}
Let $K$ be a nonempty set and $f : K \to \mathbb{R}$ a bounded
function with $m_0 = \sup f$. Let $g : K \to \mathbb{R}$ satisfy
\begin{enumerate}
\item[(i)] $g(t) \le f(t)$ for every $t \in K$;
\item[(ii)] $f(t) - g(t) \ge \varepsilon$ whenever
      $f(t) \ge m_0 - \delta$,
\end{enumerate}
where $0 < \varepsilon \le \delta$. Then
$\sup g \le m_0 - \varepsilon$.
\end{lemma}

\begin{proof}
Fix $t \in K$. If $f(t) \ge m_0 - \delta$, then (ii) gives
$g(t) \le f(t) - \varepsilon \le m_0 - \varepsilon$. If
$f(t) < m_0 - \delta$, then (i) gives $g(t) \le f(t) <
m_0 - \delta \le m_0 - \varepsilon$. Taking the supremum over
$t$ completes the proof.
\end{proof}

Lemma~\ref{lem:scalar-sup-reduction} is a general statement
about lowering the supremum of a bounded real-valued function by
exhibiting a competitor with a pointwise saving on the
super-level set $\{f \ge m_0 - \delta\}$; no sweepout, cover, or
multiplicity hypothesis enters. It applies verbatim whenever a
pointwise saving can be established on the ``where it matters''
super-level set. The Lean library ships two corollaries of
Theorem~\ref{thm:core} via this lemma: the generic-$K$ form (any
\texttt{FiniteCoverWithWitnesses}), and the one-parameter
specialization to $\unit$ that chains through
Theorem~\ref{thm:core}.

\begin{corollary}[Sup reduction, generic form]
\label{cor:core-sup}
\leanmargin{exists\_sup\_reduction\_of\_cover}{Cover.lean}
Let $K$ be a compact, nonempty topological space,
$f : K \to \mathbb{R}$ continuous with $m_0 = \sup f$, and
$0 < \varepsilon \le \delta$. Given a
\texttt{FiniteCoverWithWitnesses}$(K, f, m_0, \delta, \varepsilon)$
(Definition~\ref{def:ref}) with pieces $\{p_l\}$, there exists
$f' : K \to \mathbb{R}$ satisfying:
\begin{enumerate}
\item[\textup{(a)}] $f'(t) \le f(t)$ for every $t \in K$;
\item[\textup{(b)}] $f'(t) = f(t)$ whenever $t \notin \bigcup_l p_l$;
\item[\textup{(c)}] $\sup f' \le m_0 - \varepsilon$.
\end{enumerate}
\end{corollary}

\begin{proof}
Set $f'(t) := f(t) - \sum_{l \,:\, t \in p_l} s_l$. Since each
$s_l > 0$, (a) holds. If $t \notin \bigcup_l p_l$ the sum is
empty, giving (b). For (c), apply
Lemma~\ref{lem:scalar-sup-reduction}: (i) follows from (a);
(ii) holds because when $f(t) \ge m_0 - \delta$, (IV) gives
some $p_l \ni t$, (III) gives $s_l \ge \varepsilon$, and the
remaining at-most-one saving (by two-fold overlap (II)) is
non-negative, so $f(t) - f'(t) \ge \varepsilon$. The lemma
yields $\sup f' \le m_0 - \varepsilon$.
\end{proof}

\begin{corollary}[One-parameter sup reduction]
\label{thm:main}
\leanmargin{exists\_sup\_reduction}{SupReduction.lean}
Under the hypotheses of Theorem~\ref{thm:core}, there exist
$f' : \unit \to \mathbb{R}$ and a set $S \subseteq \unit$
containing $\NC$ such that:
\begin{enumerate}
\item[\textup{(a)}] $f'(t) \le f(t)$ for every $t \in \unit$;
\item[\textup{(b)}] $f'(t) = f(t)$ whenever $t \notin S$;
\item[\textup{(c)}] $\sup f' \le m_0 - 1/(4N)$.
\end{enumerate}
\end{corollary}

Corollary~\ref{thm:main} is Corollary~\ref{cor:core-sup} applied
to the \texttt{FiniteCoverWithWitnesses} produced by
Theorem~\ref{thm:core}, with
$(\delta, \varepsilon) = (1/N, 1/(4N))$; the modification set
$S$ is the union of the cover pieces, which contains $\NC$ by
condition~(IV) of Definition~\ref{def:ref}.

\subsection{Internal lemmas}
\label{sec:internal}

The declarations below are the sub-lemmas of
Theorem~\ref{thm:core}; proofs are sketched in
\S\ref{sec:proof}.

\begin{proposition}[Compactness of $\NC$]
\label{prop:NC-compact}
\leanmargintwo{isCompact\_nearCritical}{OneDim/InitialCover.lean}%
              {nearCritical\_nonempty}{OneDim/InitialCover.lean}
If $K$ is compact and $f$ continuous, $\NC$ is compact. If
additionally $m_0 = \sup f$ and $N > 0$, $\NC$ is nonempty.
\end{proposition}

\begin{itemize}
\item \lean{exists\_initialCover}{OneDim/CoverConstruction.lean}
: existence of an \texttt{InitialCover} via a grid plus
Lebesgue-number construction on the compact sublevel set $\NC$.
\item \lean{exists\_terminal\_refinement}{OneDim/Induction.lean}
: existence of a terminal \texttt{PartialRefinement} with
injective $\sigma$ via bounded smallest-index induction with
two-case dispatch.
\item \lean{not\_three\_overlap}{OneDim/SpacedIntervals.lean}
: at most two among any three distinct \texttt{InitialCover}
intervals overlap at a common point (parity rescue).
\item \lean{saving\_bound\_closure}{OneDim/Assembly.lean}
: the open-neighborhood saving bound in \texttt{LocalWitness}
extends to the closure via continuity of $E$.
\end{itemize}

The assembly itself, which packages the Stage A--B--C output
as a
\texttt{Finite\discretionary{-}{}{}Cover\discretionary{-}{}{}With\discretionary{-}{}{}Witnesses}
with $(\delta, \varepsilon) = (1/N, 1/(4N))$ on $\unit$, is
Theorem~\ref{thm:core} itself, whose Lean declaration
\lean{exists\_refinement}{OneDim/Assembly.lean} chains the
sub-lemmas above.

Figure~\ref{fig:arch} shows the two-tier architecture.

\begin{figure}[tb]
\centering
\begin{tikzpicture}[
  font=\small,
  every node/.style={align=center},
  input/.style={draw=black, thin, rectangle, rounded corners=3pt,
                inner sep=5pt, text width=46mm},
  stage/.style={draw=black!55, thin, rectangle, rounded corners=2pt,
                inner sep=3pt, text width=42mm, font=\scriptsize},
  stagemain/.style={draw=black, thick, rectangle, rounded corners=2pt,
                    inner sep=3pt, text width=42mm, font=\scriptsize,
                    fill=black!5},
  tierbox/.style={draw=black!40, thin, rectangle, rounded corners=4pt,
                  dashed, inner sep=8pt},
  thm/.style={draw=black, thin, rectangle, rounded corners=3pt,
              inner sep=5pt, text width=58mm},
  future/.style={draw=black!50, dashed, rectangle, rounded corners=3pt,
                 inner sep=4pt, text width=44mm, text=black!60,
                 font=\scriptsize},
  arr/.style={-Latex, thin, black!70},
  futarr/.style={-Latex, thin, black!50, dashed},
  arrlab/.style={font=\scriptsize\itshape, midway, fill=white,
                 inner sep=2pt, text=black!60},
  tierlab/.style={font=\scriptsize\bfseries\itshape, text=black!70}
]

% Input
\node[input] (input) at (0, 6.4) {%
  \texttt{LocalWitness} hypothesis \\
  {\scriptsize on the $1/N$-near-critical set of $f$}
};

% Tier 1 (left): DLT path with three stages
\node[stage] (sA) at (-4.6, 4.6) {%
  Stage A: \texttt{exists\_initialCover} \\
  Lebesgue number $\to$ skip-$2$ spacing (a)
};
\node[stage] (sB) at (-4.6, 3.3) {%
  Stage B: \texttt{exists\_terminal\_refinement} \\
  smallest-index induction $\to$ $\sigma$ injective
};
\node[stagemain] (sC) at (-4.6, 1.9) {%
  Stage C: \texttt{exists\_DLTCover} \\
  \textbf{structured output} (Stage A + B + invariants)
};

% Tier 1 grouping label
\node[tierlab] at (-4.6, 5.6)
  {Tier 1 \,---\, \texttt{CombArg.OneDim} \,---\, DLT-faithful path};

% Tier 2 (right): single direct theorem
\node[thm, text width=46mm] (sP) at (4.2, 3.3) {%
  \texttt{exists\_refinement\_partition} \\
  {\scriptsize partition by sorted endpoints}
};

% Tier 2 grouping label
\node[tierlab] at (4.2, 5.6)
  {Tier 2 \,---\, \texttt{CombArg.Scalar} \,---\, cheap path};

% Common abstract output
\node[thm] (fcw) at (0, 0.0) {%
  \texttt{FiniteCoverWithWitnesses} \\
  {\scriptsize abstract scalar output (multiplicity $\le 2$, saving $\ge 1/(4N)$)}
};

% Bookkeeping corollary
\node[thm] (sup) at (0, -2.0) {%
  \texttt{exists\_sup\_reduction} \\
  {\scriptsize sup-reducing competitor $f'$ with
   $\sup f' \le m_0 - 1/(4N)$}
};

% Future: geometric lift (dashed)
\node[future] (geom) at (-4.6, 0.0) {%
  \emph{(future)} \texttt{Geometric/} \\
  modified sweepout $\{\partial \Omega'_t\}$ via blending
};

% Arrows
\draw[arr] (input.south) -| (sA.north);
\draw[arr] (input.south) -| (sP.north);
\draw[arr] (sA.south) -- (sB.north);
\draw[arr] (sB.south) -- (sC.north);
\draw[arr] (sC.south) -- node[arrlab, sloped, above]
  {\texttt{toFinite}} (fcw.north west);
\draw[arr] (sP.south) -- (fcw.north east);
\draw[arr] (fcw.south) --
  node[arrlab, right=1pt] {bookkeeping (\texttt{Cover.lean})}
  (sup.north);
\draw[futarr] (sC.east) to[bend right=15] (geom.north);

\end{tikzpicture}
\caption{\label{fig:arch}\CombArg{} v0.4 two-tier architecture.
The witness hypothesis feeds two independent proof paths producing
the same abstract output: Tier~1 (\texttt{CombArg.OneDim}) follows
the DLT \S3.2 Step~1 algorithm in three stages culminating in the
structured \texttt{DLTCover}; Tier~2 (\texttt{CombArg.Scalar})
proves the same conclusion directly by partition by endpoints and
imports none of Tier~1's spacing / refinement / parity machinery
(the dependency-graph fact behind
Remark~\ref{rem:why-construction}). Both deliver a
\texttt{FiniteCoverWithWitnesses} that the $K$-generic bookkeeping
corollary in \texttt{Cover.lean} turns into a sup-reducing
competitor $f'$. The future geometric lift consumes the
\texttt{DLTCover} structure, not the abstract output --- the
positive-measure overlap on $J_i \cap J_{i+1}$ that DLT's
modified sweepout requires is preserved by Tier~1 and discarded
by Tier~2.}
\end{figure}

